\subsection{\texorpdfstring{第~\ref{sec:cartan-structure-equations}~节}{Cartan 结构方程}}\label{subsec:solutions-cartan-structure-equations}

\begin{description}[font=\normalfont,leftmargin=2em,labelwidth=1.5em]
\RGsolution{ex:gauss-curvature-of-sphere-and-hyperbolic-plane}{gauss-curvature-of-sphere-and-hyperbolic-plane}
\begin{enumerate}[label=(\alph*),ref=\alph*,leftmargin=*]
\item 在 $S^2\subset\mathbb{R}^3$ 上取通常的局部坐标, 它们由参数化 $\phi:(0,\pi)\times(0,2\pi)\to\mathbb{R}^3$ 给出:
\[
  \phi(\theta,\varphi)=(\sin\theta\cos\varphi,\sin\theta\sin\varphi,\cos\theta).
\]
考虑标架场 $\{X_1,X_2\}$, 其中
\[
  X_1:=\frac{\partial}{\partial\theta}\quad\text{且}\quad X_2:=\frac{\partial}{\partial\varphi}.
\]
则
\[
  \langle X_1,X_1\rangle=1,\quad \langle X_1,X_2\rangle=0
  \quad\text{且}\quad \langle X_2,X_2\rangle=\sin^2\theta,
\]
故 $E_1:=X_1$ 和 $E_2:=\frac1{\sin\theta}X_2$ 给出标准正交标架场 $\{E_1,E_2\}$, 由 $\omega^1:=\mathrm{d}\theta$ 与 $\omega^2:=\sin\theta\,\mathrm{d}\varphi$ 给出的 $\{\omega^1,\omega^2\}$ 是其对偶余标架场. 此外,
\[
  \mathrm{d}\omega^1=0\quad\text{且}\quad
  \mathrm{d}\omega^2=\cos\theta\,\mathrm{d}\theta\wedge\mathrm{d}\varphi
      =\cot\theta\,\omega^1\wedge\omega^2.
\]
第一 Cartan 结构方程
\[
  \mathrm{d}\omega^1=\omega^2\wedge\omega_2^1
   \quad\text{与}\quad \mathrm{d}\omega^2=\omega^1\wedge\omega_1^2
\]
蕴含
\[
  \mathrm{d}\omega^1(E_1,E_2)=-\omega_2^1(E_1)=\omega_1^2(E_1)
\]
和
\[
  \mathrm{d}\omega^2(E_1,E_2)=\omega_1^2(E_2).
\]
所以联络形式 $\omega_1^2$ 由
\[
  \omega_1^2=\mathrm{d}\omega^1(E_1,E_2)\omega^1+\mathrm{d}\omega^2(E_1,E_2)\omega^2
    =\cot\theta\,\omega^2=\cos\theta\,\mathrm{d}\varphi
\]
给出. 因此 $\mathrm{d}\omega_1^2=-\sin\theta\,\mathrm{d}\theta\wedge\mathrm{d}\varphi=-\omega^1\wedge\omega^2$, 由 Proposition~\ref{prop:curvature-form-and-gauss-curvature} 可知 $S^2$ 的 Gauss 曲率为 $K=1$.
\item 在 $H$ 上考虑标架场 $\{X_1,X_2\}$, 其中
\[
  X_1:=\frac{\partial}{\partial x}\quad\text{且}\quad X_2:=\frac{\partial}{\partial y}.
\]
则
\[
  \langle X_1,X_1\rangle=\frac1{y^2},\quad \langle X_1,X_2\rangle=0
  \quad\text{且}\quad \langle X_2,X_2\rangle=\frac1{y^2},
\]
故 $E_1:=yX_1$ 和 $E_2:=yX_2$ 给出标准正交标架场 $\{E_1,E_2\}$, 由 $\omega^1:=\frac1y\,\mathrm{d}x$ 与 $\omega^2:=\frac1y\,\mathrm{d}y$ 给出的 $\{\omega^1,\omega^2\}$ 是其对偶余标架场. 此外,
\[
  \mathrm{d}\omega^1=\frac1{y^2}\,\mathrm{d}x\wedge\mathrm{d}y=\omega^1\wedge\omega^2
  \quad\text{且}\quad \mathrm{d}\omega^2=0,
\]
所以联络形式 $\omega_1^2$ 由
\[
  \omega_1^2=\mathrm{d}\omega^1(E_1,E_2)\omega^1+\mathrm{d}\omega^2(E_1,E_2)\omega^2=\omega^1
\]
给出. 因此 $\mathrm{d}\omega_1^2=\frac1{y^2}\,\mathrm{d}x\wedge\mathrm{d}y=\omega^1\wedge\omega^2$, 由 Proposition~\ref{prop:curvature-form-and-gauss-curvature} 可知 $H$ 的 Gauss 曲率为 $K=-1$.
\end{enumerate}

\RGsolution{ex:ricci-and-scalar-curvature-of-radial-metric}{ricci-and-scalar-curvature-of-radial-metric}
\begin{enumerate}[label=(\alph*),ref=\alph*,leftmargin=*]
\item 一组标准正交余标架由下式给出:
\[
  \omega^r=A(r)\,\mathrm{d}r,\quad\omega^\theta=r\,\mathrm{d}\theta
  \quad\text{且}\quad \omega^\varphi=r\sin\theta\,\mathrm{d}\varphi.
\]
由 Cartan 结构方程得
\begin{align*}
  \mathrm{d}\omega^r&=0=\omega^\theta\wedge\omega_\theta^r+\omega^\varphi\wedge\omega_\varphi^r;\\
  \mathrm{d}\omega^\theta&=\mathrm{d}r\wedge\mathrm{d}\theta
   =\omega^r\wedge\omega_r^\theta+\omega^\varphi\wedge\omega_\varphi^\theta;\\
  \mathrm{d}\omega^\varphi&=\sin\theta\,\mathrm{d}r\wedge\mathrm{d}\varphi+
        r\cos\theta\,\mathrm{d}\theta\wedge\mathrm{d}\varphi\\
   &=\omega^r\wedge\omega_r^\varphi+\omega^\theta\wedge\omega_\theta^\varphi,
\end{align*}
由此容易猜到
\begin{align*}
  \omega_r^\theta&=-\omega_\theta^r=\frac1A\,\mathrm{d}\theta;\\
  \omega_r^\varphi&=-\omega_\varphi^r=\frac{\sin\theta}{A}\,\mathrm{d}\varphi;\\
  \omega_\theta^\varphi&=-\omega_\varphi^\theta=\cos\theta\,\mathrm{d}\varphi.
\end{align*}
曲率形式由下式给出:
\begin{align*}
  \Omega_r^\theta
  &=\mathrm{d}\omega_r^\theta-\omega_r^\varphi\wedge\omega_\varphi^\theta
    =-\frac{A'}{A^2}\,\mathrm{d}r\wedge\mathrm{d}\theta=-\frac{A'}{rA^3}\,\omega^r\wedge\omega^\theta;\\
  \Omega_r^\varphi
  &=\mathrm{d}\omega_r^\varphi-\omega_r^\theta\wedge\omega_\theta^\varphi
    =-\frac{A'\sin\theta}{A^2}\,\mathrm{d}r\wedge\mathrm{d}\varphi=-\frac{A'}{rA^3}\,\omega^r\wedge\omega^\varphi;\\
  \Omega_\theta^\varphi
  &=\mathrm{d}\omega_\theta^\varphi-\omega_\theta^r\wedge\omega_r^\varphi
    =\left(\frac{\sin\theta}{A^2}-\sin\theta\right)\mathrm{d}\theta\wedge\mathrm{d}\varphi\\
  &=\frac1{r^2}\left(\frac1{A^2}-1\right)\omega^\theta\wedge\omega^\varphi,
\end{align*}
所以曲率张量在这组标准正交标架下的非零分量是
\begin{align*}
  R_{r\theta r}{}^\theta&=\Omega_r^\theta(E_r,E_\theta)=-\frac{A'}{rA^3};\\
  R_{r\varphi r}{}^\varphi&=\Omega_r^\varphi(E_r,E_\varphi)=-\frac{A'}{rA^3};\\
  R_{\theta\varphi\theta}{}^\varphi&=\Omega_\theta^\varphi(E_\theta,E_\varphi)
    =\frac1{r^2}\left(\frac1{A^2}-1\right)
\end{align*}
(以及由对称性与这些分量相联系的分量). 因此 Ricci 张量在这组标准正交标架下的分量为
\begin{align*}
  R_{rr}&=R_{\theta rr}{}^\theta+R_{\varphi rr}{}^\varphi=\frac{2A'}{rA^3};\\
  R_{\theta\theta}&=R_{r\theta\theta}{}^r+R_{\varphi\theta\theta}{}^\varphi
    =\frac{A'}{rA^3}-\frac1{r^2}\left(\frac1{A^2}-1\right);\\
  R_{\varphi\varphi}&=R_{r\varphi\varphi}{}^r+R_{\theta\varphi\varphi}{}^\theta
    =\frac{A'}{rA^3}-\frac1{r^2}\left(\frac1{A^2}-1\right).
\end{align*}
于是标量曲率为
\[
  S=R_{rr}+R_{\theta\theta}+R_{\varphi\varphi}
      =\frac{4A'}{rA^3}-\frac2{r^2}\left(\frac1{A^2}-1\right).
\]
\item 在这种情形下, 有 $R_{rr}=R_{\theta\theta}=R_{\varphi\varphi}=2$, 即 $\mathit{Ric}=2g$, 因而 $S=6$.
\item 在这种情形下, 有 $R_{rr}=R_{\theta\theta}=R_{\varphi\varphi}=-2$, 即 $\mathit{Ric}=-2g$, 因而 $S=-6$.
\item 对常数 $S$, 要解微分方程
\[
  \frac{4A'}{rA^3}-\frac2{r^2}\left(\frac1{A^2}-1\right)=S.
\]
这个方程可以写成
\[
  -\frac{2A'r}{A^3}+\frac1{A^2}=1-\frac{Sr^2}{2}
   \iff \left(\frac r{A^2}\right)'=1-\frac{Sr^2}{2},
\]
立即可得解
\[
  A(r)=\left(1-\frac{Sr^2}{6}+\frac Cr\right)^{-\frac12}
\]
(其中 $C\in\mathbb{R}$ 为积分常数).
\end{enumerate}

\RGsolution{ex:gauss-curvature-by-holonomy}{gauss-curvature-by-holonomy}
定义正定向的标准正交标架场 $\{E_1,E_2\}$ 和 $\{F_1,F_2\}$, 使得 $\{E_1,E_2\}$ 在 $\overline D$ 上有定义且 $F_1=X$. 由 Proposition~\ref{prop:connection-form-under-oriented-frame-change} 得
\[
  \Delta\theta=\int_{\partial D}\sigma=\int_{\partial D}\overline\omega_1^2-\omega_1^2.
\]
因为 $X$ 沿 $\partial D$ 平行移动, 有
\[
  \nabla_{\dot c}F_1=0\iff\overline\omega_1^2(\dot c)=0\iff c^*\overline\omega_1^2=0.
\]
因此
\[
  \Delta\theta=-\int_{\partial D}\omega_1^2=-\int_D\mathrm{d}\omega_1^2
   =\int_DK\,\omega^1\wedge\omega^2=\int_DK,
\]
其中用了 Stokes 定理. $K(p)$ 的极限公式可以用标准论证得出.

\RGsolution{ex:geodesic-curvature-as-angle-derivative}{geodesic-curvature-as-angle-derivative}
考虑正定向的两个标准正交标架场 $\{E_1,E_2\}$ 和 $\{F_1,F_2\}$, 使得 $E_1=\frac X{\lVert X\rVert}$ 且 $F_1=\dot c$, 则有
\[
  \frac{\mathrm{d}\theta}{\mathrm{d}s}(s)=\mathrm{d}\theta(\dot c(s))
  =\overline\omega_1^2(\dot c(s))-\omega_1^2(\dot c(s))
  =\overline\omega_1^2(F_1)=k_g(s),
\]
其中使用了
\[
  \nabla_{\dot c}E_1=0\iff\omega_1^2(\dot c(s))=0.
\]
\end{description}
